\documentclass[10pt,a4paper]{article}

\usepackage[utf8]{inputenc}
\usepackage[T1]{fontenc}
\usepackage[french]{babel}

\usepackage{styles/couleurs}
\usepackage{styles/mise_en_page}
\usepackage{styles/boites}
\usepackage{styles/maths}
\usepackage{styles/figures}

\usepackage{tikz}
\usetikzlibrary{
  calc,
  angles,
  quotes,
  decorations.markings
}
\usepackage{tabularx}
\usepackage{array}
\usepackage{amssymb}

% geometry est déjà chargé par mise_en_page
\geometry{
  a4paper,
  landscape,
  top=0.55cm,
  bottom=0.55cm,
  left=0.7cm,
  right=0.7cm
}

\pagestyle{empty}
\setlength{\parindent}{0pt}
\setlength{\parskip}{0pt}

% =========================================================
% TIKZ
% =========================================================

\tikzset{
  point/.style={
    circle,
    fill,
    inner sep=1.4pt
  },
  trait/.style={
    thick
  },
  codeun/.style={
    postaction={decorate},
    decoration={
      markings,
      mark=at position 0.5 with {
        \draw[thick] (0,-3pt)--(0,3pt);
      }
    }
  },
  codedeux/.style={
    postaction={decorate},
    decoration={
      markings,
      mark=at position 0.5 with {
        \draw[thick] (-1.5pt,-3pt)--(-1.5pt,3pt);
        \draw[thick] (1.5pt,-3pt)--(1.5pt,3pt);
      }
    }
  }
}

% =========================================================
% MISE EN PAGE
% =========================================================

\newcommand{\titreRDV}[2]{%
{\sffamily\LARGE\bfseries\color{TitreC}#1}
\hfill
{\sffamily\small #2}

\vspace{0.16cm}
}

\newcommand{\flashcalcul}[1]{%
\begin{tcolorbox}[
  enhanced,
  colback=SavoirBg,
  colframe=SavoirC!55,
  boxrule=0.6pt,
  arc=1.5mm,
  left=3mm,
  right=3mm,
  top=1.5mm,
  bottom=1.5mm
]
\textbf{\sffamily Calcul réfléchi}
\hfill
\small Réponse seule. Pas de calculatrice.

\vspace{0.15cm}

#1
\end{tcolorbox}
}

\newcommand{\flashgeo}[1]{%
\begin{tcolorbox}[
  enhanced,
  colback=SavoirBg,
  colframe=SavoirC!55,
  boxrule=0.6pt,
  arc=1.5mm,
  left=3mm,
  right=3mm,
  top=1.5mm,
  bottom=1.5mm
]
\textbf{\sffamily Flash géométrie}

\vspace{0.1cm}

#1
\end{tcolorbox}
}

\newcommand{\calcquatre}[4]{%
\begin{tabularx}{\linewidth}{
>{\centering\arraybackslash}X
>{\centering\arraybackslash}X
>{\centering\arraybackslash}X
>{\centering\arraybackslash}X
}
#1 & #2 & #3 & #4 \\[0.25cm]
\rule{3.7cm}{0.4pt} &
\rule{3.7cm}{0.4pt} &
\rule{3.7cm}{0.4pt} &
\rule{3.7cm}{0.4pt}
\end{tabularx}
}

\newcommand{\grandeactivite}[2]{%
\begin{tcolorbox}[
  enhanced,
  colback=white,
  colframe=NeutreC!55,
  boxrule=0.6pt,
  arc=1.5mm,
  left=4mm,
  right=4mm,
  top=2.5mm,
  bottom=2.5mm,
  title={\sffamily\bfseries\color{TitreC}#1},
  coltitle=TitreC,
  colbacktitle=NeutreBg
]
#2
\end{tcolorbox}
}

% =========================================================
% DOCUMENT
% =========================================================

\begin{document}

% =========================================================
% RENDEZ-VOUS 1
% FIGURE TÉLÉPHONÉE
%
% Cette page est à imprimer puis à couper en deux.
% Une moitié de la classe reçoit A, l'autre reçoit B.
%
% Les dessins sources sont exacts.
% Le croquis produit par l'élève est à main levée
% et n'est pas à l'échelle.
% =========================================================

\titreRDV
{Rendez-vous 1 -- Transmettre une figure}
{Modèles A et B -- à couper}

\flashcalcul{
\calcquatre
{\(57+29\)}
{\(403-197\)}
{\(39\times7\)}
{\(420\div5\)}
}

\vspace{0.2cm}

\begin{minipage}[t]{0.485\linewidth}

\begin{tcolorbox}[
  enhanced,
  height=14.7cm,
  colback=white,
  colframe=NeutreC!55,
  boxrule=0.7pt,
  arc=1.5mm,
  left=4mm,
  right=4mm,
  top=2mm,
  bottom=2mm,
  title={\sffamily\bfseries\color{TitreC}Modèle A},
  coltitle=TitreC,
  colbacktitle=NeutreBg
]

\begin{center}
\textbf{Le dessin ci-dessous est exact.}

\smallskip

Tu peux utiliser tes instruments pour l'analyser.
\end{center}

\vspace{0.2cm}

\begin{center}
\begin{tikzpicture}[scale=0.83]

% Rectangle exact 8 x 5
\coordinate (A) at (0,0);
\coordinate (B) at (8,0);
\coordinate (C) at (8,5);
\coordinate (D) at (0,5);

\coordinate (M) at (4,0);
\coordinate (N) at (4,5);
\coordinate (O) at (4,2.5);

\draw[trait] (A)--(B)--(C)--(D)--cycle;
\draw[trait] (A)--(C);
\draw[trait] (B)--(D);
\draw[trait] (M)--(N);

\node[below left] at (A) {$A$};
\node[below right] at (B) {$B$};
\node[above right] at (C) {$C$};
\node[above left] at (D) {$D$};

\node[below] at (M) {$M$};
\node[above] at (N) {$N$};
\node[right] at (O) {$O$};

\end{tikzpicture}
\end{center}

\vfill

\textbf{Ta mission}

\medskip

Sur une feuille à part, réalise un \textbf{croquis à main levée}
qui permette à un camarade de reconstruire ce dessin.

\medskip

Ton croquis \textbf{n'a pas besoin d'être à l'échelle}.

\medskip

Tu peux écrire et coder toutes les informations
que tu juges nécessaires.

\medskip

\textbf{Ne montre pas le modèle à ton camarade.}

\vspace{0.3cm}

Quand ton croquis est prêt, rends ce modèle au professeur
et donne seulement ton croquis à un élève qui avait le modèle B.

\medskip

S'il lui manque une information,
il doit écrire sa question au lieu de deviner.

\end{tcolorbox}

\end{minipage}
\hfill
\begin{minipage}[t]{0.485\linewidth}

\begin{tcolorbox}[
  enhanced,
  height=14.7cm,
  colback=white,
  colframe=NeutreC!55,
  boxrule=0.7pt,
  arc=1.5mm,
  left=4mm,
  right=4mm,
  top=2mm,
  bottom=2mm,
  title={\sffamily\bfseries\color{TitreC}Modèle B},
  coltitle=TitreC,
  colbacktitle=NeutreBg
]

\begin{center}
\textbf{Le dessin ci-dessous est exact.}

\smallskip

Tu peux utiliser tes instruments pour l'analyser.
\end{center}

\vspace{0.2cm}

\begin{center}
\begin{tikzpicture}[scale=1]

% Losange exact de côté 4
% A(0,0), B(4,0), C(6,3.464), D(2,3.464)

\coordinate (A) at (0,0);
\coordinate (B) at (4,0);
\coordinate (C) at (6,3.464);
\coordinate (D) at (2,3.464);

\coordinate (O) at (3,1.732);

\draw[trait] (A)--(B)--(C)--(D)--cycle;
\draw[trait] (A)--(C);
\draw[trait] (B)--(D);

\node[below left] at (A) {$A$};
\node[below right] at (B) {$B$};
\node[above right] at (C) {$C$};
\node[above left] at (D) {$D$};

\node[right] at (O) {$O$};

\end{tikzpicture}
\end{center}

\vfill

\textbf{Ta mission}

\medskip

Sur une feuille à part, réalise un \textbf{croquis à main levée}
qui permette à un camarade de reconstruire ce dessin.

\medskip

Ton croquis \textbf{n'a pas besoin d'être à l'échelle}.

\medskip

Tu peux écrire et coder toutes les informations
que tu juges nécessaires.

\medskip

\textbf{Ne montre pas le modèle à ton camarade.}

\vspace{0.3cm}

Quand ton croquis est prêt, rends ce modèle au professeur
et donne seulement ton croquis à un élève qui avait le modèle A.

\medskip

S'il lui manque une information,
il doit écrire sa question au lieu de deviner.

\end{tcolorbox}

\end{minipage}

% =========================================================
% PROF -- RENDEZ-VOUS 1
%
% Calcul :
%
% 57+29 = 86
% -> 57+30-1
%
% 403-197 = 206
% -> 403-200+3
%
% 39x7 = 273
% -> 40x7-7
%
% 420÷5 = 84
% -> 840÷10
%
%
% FIGURE A
%
% Rectangle 8 x 5 exactement.
% M et N sont les milieux de [AB] et [DC].
% O est l'intersection des diagonales,
% et appartient également à [MN].
%
% Les élèves peuvent mesurer puisque le modèle est exact.
%
%
% FIGURE B
%
% Losange exact de côté 4.
% Les diagonales sont tracées.
%
% Intérêt didactique :
% "losange de côté 4 cm" ne suffit pas à déterminer
% un losange unique.
%
% Une mesure supplémentaire pertinente
% (diagonale, angle...) devient nécessaire.
%
%
% ORGANISATION
%
% 1. Distribuer A à une moitié, B à l'autre.
% 2. Analyse individuelle avec instruments.
% 3. Croquis sur feuille séparée.
% 4. Récupérer les modèles exacts.
% 5. Échanger uniquement les croquis A <-> B.
% 6. Reconstruction.
% 7. Si une donnée manque : écrire une question.
%
%
% MISE EN COMMUN
%
% Choisir peu de productions :
% - une reconstruction réussie ;
% - une reconstruction possible mais différente ;
% - un destinataire ayant dû demander une information.
%
% Questions :
%
% "Qu'est-ce qui était vrai sur le dessin exact
% mais qui n'a pas été transmis ?"
%
% "Qu'est-ce qu'un croquis doit conserver ?"
%
% Ne pas institutionnaliser toute la liste avant le travail.
% =========================================================

\newpage

% =========================================================
% RENDEZ-VOUS 2
% DESSIN / INFORMATIONS
% =========================================================

\titreRDV
{Rendez-vous 2 -- Le dessin n'est pas la figure}
{Deux tâches courtes}

\flashcalcul{
\calcquatre
{\(25\times44\)}
{\(64\times5\)}
{\(125\times16\)}
{\(98\times23\)}
}

\vspace{0.15cm}

\grandeactivite{A. Le même croquis, des informations différentes}{

Les trois croquis sont identiques.

\medskip

\textbf{Ils ne sont pas à l'échelle.}

\vspace{0.15cm}

\begin{minipage}[t]{0.32\linewidth}
\centering

\textbf{1}

\begin{tikzpicture}[scale=0.6]

\coordinate (A) at (0,0);
\coordinate (B) at (4.6,0.35);
\coordinate (C) at (5.4,2.2);
\coordinate (D) at (0.8,1.85);

\draw[trait] (A)--(B)--(C)--(D)--cycle;

\node[below left] at (A) {$A$};
\node[below right] at (B) {$B$};
\node[above right] at (C) {$C$};
\node[above left] at (D) {$D$};

\end{tikzpicture}

Aucune autre information.

\medskip

Nature la plus précise :

\vspace{0.25cm}

\rule{5cm}{0.4pt}

\end{minipage}
\hfill
\begin{minipage}[t]{0.32\linewidth}
\centering

\textbf{2}

\begin{tikzpicture}[scale=0.6]

\coordinate (A) at (0,0);
\coordinate (B) at (4.6,0.35);
\coordinate (C) at (5.4,2.2);
\coordinate (D) at (0.8,1.85);

\draw[trait] (A)--(B)--(C)--(D)--cycle;

\node[below left] at (A) {$A$};
\node[below right] at (B) {$B$};
\node[above right] at (C) {$C$};
\node[above left] at (D) {$D$};

\end{tikzpicture}

\[
(AB)\parallel(CD)
\qquad
(AD)\parallel(BC)
\]

Nature la plus précise :

\vspace{0.25cm}

\rule{5cm}{0.4pt}

\end{minipage}
\hfill
\begin{minipage}[t]{0.32\linewidth}
\centering

\textbf{3}

\begin{tikzpicture}[scale=0.6]

\coordinate (A) at (0,0);
\coordinate (B) at (4.6,0.35);
\coordinate (C) at (5.4,2.2);
\coordinate (D) at (0.8,1.85);

\draw[trait] (A)--(B)--(C)--(D)--cycle;

\node[below left] at (A) {$A$};
\node[below right] at (B) {$B$};
\node[above right] at (C) {$C$};
\node[above left] at (D) {$D$};

\end{tikzpicture}

\[
AB=BC=CD=DA
\]

Nature la plus précise :

\vspace{0.25cm}

\rule{5cm}{0.4pt}

\end{minipage}

% PROF :
% 1 : on peut seulement garantir "quadrilatère".
% 2 : parallélogramme.
% 3 : losange.
}

\vspace{0.15cm}

\grandeactivite{B. Des croquis différents, les mêmes informations}{

Dans les trois cas :

\[
AB=AC.
\]

Les croquis ne sont pas à l'échelle.

\medskip

\begin{minipage}[c]{0.26\linewidth}

Écris une affirmation qui est
\textbf{certaine dans les trois cas}.

\vspace{0.5cm}

\rule{\linewidth}{0.4pt}

\vspace{0.5cm}

\rule{\linewidth}{0.4pt}

\end{minipage}
\hfill
\begin{minipage}[c]{0.70\linewidth}

\begin{center}
\begin{tikzpicture}[scale=0.72]

% triangle 1
\begin{scope}[shift={(0,0)}]
\coordinate (A) at (1.7,2.5);
\coordinate (B) at (0,0);
\coordinate (C) at (3.4,0);

\draw[trait] (A)--(B);
\draw[trait] (A)--(C);
\draw[trait] (B)--(C);

\node[above] at (A) {$A$};
\node[below] at (B) {$B$};
\node[below] at (C) {$C$};
\end{scope}

% triangle 2
\begin{scope}[shift={(5.2,0)}]
\coordinate (A) at (0.6,2.3);
\coordinate (B) at (0,0);
\coordinate (C) at (4.4,0.5);

\draw[trait] (A)--(B);
\draw[trait] (A)--(C);
\draw[trait] (B)--(C);

\node[above] at (A) {$A$};
\node[below] at (B) {$B$};
\node[right] at (C) {$C$};
\end{scope}

% triangle 3
\begin{scope}[shift={(11.5,0)}]
\coordinate (A) at (3.4,1.0);
\coordinate (B) at (0,0);
\coordinate (C) at (0.5,3.2);

\draw[trait] (A)--(B);
\draw[trait] (A)--(C);
\draw[trait] (B)--(C);

\node[right] at (A) {$A$};
\node[below] at (B) {$B$};
\node[above] at (C) {$C$};
\end{scope}

\end{tikzpicture}
\end{center}

\end{minipage}

% =========================================================
% PROF -- RENDEZ-VOUS 2
%
% Calcul :
%
% 25x44 = 1100
% 44 = 4x11 -> 25x4x11
%
% 64x5 = 320
% x10 puis moitié
%
% 125x16 = 2000
% 125x8=1000 puis doubler
%
% 98x23 = 2254
% 100x23 - 2x23
%
% Partie B :
% ABC est isocèle en A.
%
% On peut aussi accepter :
% angles ABC et BCA de même mesure,
% si cette propriété est connue.
% =========================================================
}

\newpage

% =========================================================
% RENDEZ-VOUS 3
% INFORMATION SUFFISANTE / INSUFFISANTE
% =========================================================

\titreRDV
{Rendez-vous 3 -- A-t-on assez d'informations ?}
{Construction}

\flashgeo{

\begin{minipage}[c]{0.27\linewidth}

\begin{center}
\begin{tikzpicture}[scale=0.7]

\node[point,label=below:$A$] at (0,0) {};
\node[point,label=below:$B$] at (2.2,0.05) {};
\node[point,label=below:$C$] at (4.5,0) {};

\end{tikzpicture}
\end{center}

\end{minipage}
\hfill
\begin{minipage}[c]{0.69\linewidth}

\textbf{Affirmation :}
\(A\), \(B\) et \(C\) sont alignés.

\medskip

$\square$ Certain
\qquad
$\square$ Faux
\qquad
$\square$ On ne peut pas savoir

\end{minipage}

% PROF :
% On ne peut pas savoir.
}

\vspace{0.15cm}

\grandeactivite{Deux figures possibles}{

On veut construire un quadrilatère \(ABCD\) tel que :

\[
AB=CD=4\text{ cm}
\qquad\text{et}\qquad
BC=DA=3\text{ cm}.
\]

\vspace{0.15cm}

\textbf{Construis deux quadrilatères différents}
qui respectent ces informations
et qui ne sont pas superposables.

\vspace{0.25cm}

\begin{minipage}[t]{0.48\linewidth}

\begin{center}
\textbf{Construction 1}
\end{center}

\begin{tcolorbox}[
  colback=white,
  colframe=NeutreC!30,
  boxrule=0.4pt,
  height=7.5cm
]
\end{tcolorbox}

\end{minipage}
\hfill
\begin{minipage}[t]{0.48\linewidth}

\begin{center}
\textbf{Construction 2}
\end{center}

\begin{tcolorbox}[
  colback=white,
  colframe=NeutreC!30,
  boxrule=0.4pt,
  height=7.5cm
]
\end{tcolorbox}

\end{minipage}

\vspace{0.35cm}

\textbf{Ajoute une seule information pour être certain que \(ABCD\) est un rectangle :}

\vspace{0.25cm}

\rule{15cm}{0.4pt}

% =========================================================
% PROF -- RENDEZ-VOUS 3
%
% Les données imposent des côtés opposés de même longueur :
% le quadrilatère est un parallélogramme.
%
% On peut construire par exemple :
% - un rectangle 4 x 3 ;
% - un parallélogramme oblique de côtés 4 et 3.
%
% Une information suffisante :
% (AB) ⟂ (BC)
%
% Si la propriété du parallélogramme n'est pas disponible,
% accepter une information plus forte comme :
% "ABCD possède quatre angles droits".
%
% Relances :
%
% 1. "Les longueurs fixent-elles les angles ?"
%
% 2. "Peux-tu garder les quatre longueurs
%     tout en inclinant la figure ?"
% =========================================================
}

\newpage

% =========================================================
% RENDEZ-VOUS 4
% RESTAURATION
% =========================================================

\titreRDV
{Rendez-vous 4 -- Les propriétés deviennent des outils}
{Restauration}

\flashcalcul{
\calcquatre
{\(6\times\Box=4{,}8\)}
{\(\Box\div7=9\)}
{\(42\div\Box=6\)}
{\(20\%\text{ de }\Box=14\)}
}

\vspace{0.15cm}

\grandeactivite{Le losange effacé}{

\begin{minipage}[t]{0.30\linewidth}

On sait que \(ABCD\) était un \textbf{losange}.

Trois sommets ont été conservés.

\medskip

\textbf{Restaure la figure.}

\medskip

Instruments autorisés :

\[
\text{règle non graduée}
\qquad
\text{compas}
\]

\vspace{0.5cm}

Quand tu penses avoir terminé,
demande le calque de validation.

\vspace{0.7cm}

\textbf{Après validation :}

Quelle propriété du losange
t'a permis de placer \(D\) ?

\vspace{0.35cm}

\rule{\linewidth}{0.4pt}

\vspace{0.4cm}

\rule{\linewidth}{0.4pt}

\end{minipage}
\hfill
\begin{minipage}[t]{0.66\linewidth}

\begin{center}
\begin{tikzpicture}[scale=1.6]

% AB = BC = 2.5

\coordinate (B) at (0,0);
\coordinate (A) at (-2.5,0);
\coordinate (C) at (1.5,2);

\node[point,label=below:$A$] at (A) {};
\node[point,label=below:$B$] at (B) {};
\node[point,label=above right:$C$] at (C) {};

\draw[trait] (A)--(B)--(C);

\end{tikzpicture}
\end{center}

\vspace{0.4cm}

\begin{tcolorbox}[
  colback=white,
  colframe=NeutreC!20,
  boxrule=0.4pt,
  height=5.8cm
]
\textbf{Traces de recherche}
\end{tcolorbox}

\end{minipage}

% =========================================================
% PROF -- RENDEZ-VOUS 4
%
% Calcul :
%
% 6 x ? = 4,8        -> 0,8
% ? ÷ 7 = 9          -> 63
% 42 ÷ ? = 6         -> 7
% 20 % de ? = 14     -> 70
%
%
% Restauration :
%
% D doit vérifier :
% AD = AB
% CD = CB
%
% Construction possible :
% cercle de centre A et rayon AB
% cercle de centre C et rayon CB
%
% Les deux intersections sont B et D.
%
% Relances graduées :
%
% 1. "Que sais-tu des côtés d'un losange ?"
%
% 2. "À quelle distance de A doit être D ?"
%
% 3. "À quelle distance de C doit être D ?"
%
% Ne pas donner la procédure avant la recherche.
% =========================================================
}

\newpage

% =========================================================
% RENDEZ-VOUS 5
% DONNÉ / DÉDUIT / NON GARANTI
% =========================================================

\titreRDV
{Rendez-vous 5 -- Donné, déduit, non garanti}
{Deux tâches intermédiaires}

\flashcalcul{
\calcquatre
{\(0{,}5\times74\)}
{\(50\%\text{ de }46\)}
{\(\dfrac34\text{ de }44\)}
{\(0{,}25\times68\)}
}

\vspace{0.15cm}

\grandeactivite{A. Quel est le statut de chaque affirmation ?}{

\begin{minipage}[c]{0.32\linewidth}

On sait que \(ABCD\) est un rectangle.

\(O\) est le point d'intersection de ses diagonales.

\medskip

\textbf{Le dessin est un croquis.}

\begin{center}
\begin{tikzpicture}[scale=0.75]

\coordinate (A) at (0,0);
\coordinate (B) at (4.8,0.4);
\coordinate (C) at (4.45,2.5);
\coordinate (D) at (-0.35,2.1);
\coordinate (O) at (2.22,1.25);

\draw[trait] (A)--(B)--(C)--(D)--cycle;
\draw[trait] (A)--(C);
\draw[trait] (B)--(D);

\node[below left] at (A) {$A$};
\node[below right] at (B) {$B$};
\node[above right] at (C) {$C$};
\node[above left] at (D) {$D$};
\node[above] at (O) {$O$};

\end{tikzpicture}
\end{center}

\end{minipage}
\hfill
\begin{minipage}[c]{0.64\linewidth}

Pour chaque affirmation, écris :

\[
\boxed{\text{Donnée}}
\qquad
\boxed{\text{Déduite}}
\qquad
\boxed{\text{Non garantie}}
\]

\renewcommand{\arraystretch}{1.4}

\begin{tabularx}{\linewidth}{|X|p{4cm}|}
\hline
\(ABCD\) est un rectangle. & \\ \hline
\((AB)\perp(BC)\) & \\ \hline
\(AC=BD\) & \\ \hline
\(AO=OC\) & \\ \hline
\((AC)\perp(BD)\) & \\ \hline
\(AB=BC\) & \\ \hline
\end{tabularx}

\end{minipage}

% PROF :
% donnée
% déduite
% déduite
% déduite
% non garantie
% non garantie
}

\vspace{0.15cm}

\grandeactivite{B. Organiser un raisonnement}{

Remets les trois informations dans un ordre logique.

\vspace{0.35cm}

\begin{center}
\begin{tikzpicture}

\node[
draw,
rounded corners,
minimum width=6cm,
minimum height=1.5cm
] at (0,0)
{\(AB=AC\)};

\node[
draw,
rounded corners,
minimum width=6cm,
minimum height=1.5cm
] at (8,0)
{\(ABC\) est isocèle en \(A\)};

\node[
draw,
rounded corners,
minimum width=7cm,
minimum height=1.5cm
] at (16,0)
{\(\widehat{ABC}=\widehat{BCA}\)};

\end{tikzpicture}
\end{center}

\vspace{0.45cm}

\textbf{Qu'est-ce qui permet de passer d'une affirmation à la suivante ?}

\vspace{0.25cm}

\rule{\linewidth}{0.4pt}

% =========================================================
% PROF -- RENDEZ-VOUS 5
%
% Calcul :
%
% 0,5 x 74 = 37
% 50 % de 46 = 23
% 3/4 de 44 = 33
% 0,25 x 68 = 17
%
% Graphe :
%
% AB=AC
% -> définition / reconnaissance d'un triangle isocèle
% -> ABC isocèle en A
% -> propriété des angles à la base
% -> ABC=BCA
%
% Pas de rédaction formelle imposée.
% =========================================================
}

\newpage

% =========================================================
% RENDEZ-VOUS 6
% CONTRE-EXEMPLE
% =========================================================

\titreRDV
{Rendez-vous 6 -- Vrai parce que... faux parce que...}
{Travail en groupe}

\flashcalcul{
\calcquatre
{\(18\times6+18\times4\)}
{\(25\times63+25\times37\)}
{\(32\times2{,}5\)}
{\(125\times32\)}
}

\vspace{0.15cm}

\grandeactivite{Convaincre}{

Choisissez \textbf{deux affirmations}.

\medskip

Si vous pensez qu'elle est fausse,
\textbf{fabriquez un contre-exemple}.

Si vous pensez qu'elle est vraie,
\textbf{expliquez pourquoi elle doit toujours être vraie}.

\vspace{0.3cm}

\begin{minipage}[t]{0.48\linewidth}

\begin{tcolorbox}[
  colback=white,
  colframe=NeutreC!45,
  boxrule=0.5pt,
  height=5.3cm
]
\textbf{A}

\vspace{0.3cm}

Si un quadrilatère a quatre côtés de même longueur,
alors c'est un carré.

\vfill

\textbf{Recherche :}
\end{tcolorbox}

\vspace{0.2cm}

\begin{tcolorbox}[
  colback=white,
  colframe=NeutreC!45,
  boxrule=0.5pt,
  height=5.3cm
]
\textbf{C}

\vspace{0.3cm}

Si un triangle a deux côtés de même longueur,
alors il a deux angles de même mesure.

\vfill

\textbf{Recherche :}
\end{tcolorbox}

\end{minipage}
\hfill
\begin{minipage}[t]{0.48\linewidth}

\begin{tcolorbox}[
  colback=white,
  colframe=NeutreC!45,
  boxrule=0.5pt,
  height=5.3cm
]
\textbf{B}

\vspace{0.3cm}

Si un quadrilatère a quatre angles droits,
alors c'est un carré.

\vfill

\textbf{Recherche :}
\end{tcolorbox}

\vspace{0.2cm}

\begin{tcolorbox}[
  colback=white,
  colframe=NeutreC!45,
  boxrule=0.5pt,
  height=5.3cm
]
\textbf{D}

\vspace{0.3cm}

Si deux droites sont perpendiculaires
à une même droite,
alors elles sont parallèles.

\vfill

\textbf{Recherche :}
\end{tcolorbox}

\end{minipage}

% =========================================================
% PROF -- RENDEZ-VOUS 6
%
% Calcul :
%
% 18x6 + 18x4
% = 18x10 = 180
%
% 25x63 + 25x37
% = 25x100 = 2500
%
% 32x2,5
% = 32x5/2 = 80
%
% 125x32
% = 125x8x4
% = 1000x4
% = 4000
%
%
% A faux : losange non carré.
% B faux : rectangle non carré.
% C vrai.
% D vrai.
%
% Point à faire émerger :
%
% Un seul contre-exemple suffit pour rendre
% une affirmation générale fausse.
%
% Plusieurs exemples ne prouvent pas
% qu'une affirmation générale est vraie.
% =========================================================
}

\newpage

% =========================================================
% RENDEZ-VOUS 7
% SITUATION DENSE
% =========================================================

\titreRDV
{Rendez-vous 7 -- Exactement ou à peu près ?}
{Situation de recherche}

\flashgeo{

\begin{minipage}[c]{0.45\linewidth}

\begin{center}
\begin{tikzpicture}[scale=0.65]

\coordinate (A) at (0,0);
\coordinate (B) at (3,0);
\coordinate (C) at (6,0);
\coordinate (D) at (1.3,2.3);

\draw[trait,codeun] (A)--(D);
\draw[trait,codeun] (D)--(B);
\draw[trait,codeun] (A)--(B);
\draw[trait] (D)--(C);
\draw[trait,codeun] (B)--(C);

\node[below] at (A) {$A$};
\node[below] at (B) {$B$};
\node[below] at (C) {$C$};
\node[above] at (D) {$D$};

\end{tikzpicture}
\end{center}

\end{minipage}
\hfill
\begin{minipage}[c]{0.51\linewidth}

\textbf{Ce dessin est un croquis.}

Les quatre segments marqués ont la même longueur.

\medskip

Écris une question mathématique
à laquelle on peut répondre exactement.

\vspace{0.3cm}

\rule{\linewidth}{0.4pt}

\end{minipage}
}

\vspace{0.15cm}

\grandeactivite{Quelle est la nature exacte du triangle \(ACD\) ?}{

\begin{minipage}[c]{0.34\linewidth}

\textbf{Le dessin est un croquis.}

\medskip

Les points \(A\), \(B\) et \(C\) sont alignés.

\medskip

Les quatre segments marqués ont la même longueur.

\vspace{0.7cm}

{\large
\textbf{Quelle est la nature exacte du triangle \(ACD\) ?}
}

\vspace{0.8cm}

Tu peux écrire, coder et ajouter
des mesures sur le croquis.

\vspace{0.8cm}

\textbf{Conclusion :}

\vspace{0.35cm}

\rule{8cm}{0.5pt}

\end{minipage}
\hfill
\begin{minipage}[c]{0.62\linewidth}

\begin{center}
\begin{tikzpicture}[scale=1.45]

% Croquis volontairement non conforme.

\coordinate (A) at (0,0);
\coordinate (B) at (3,0);
\coordinate (C) at (6,0);
\coordinate (D) at (1.15,2.45);

\draw[very thick,codeun] (A)--(D);
\draw[very thick,codeun] (D)--(B);
\draw[very thick] (D)--(C);
\draw[very thick,codeun] (A)--(B);
\draw[very thick,codeun] (B)--(C);

\node[below] at (A) {$A$};
\node[below] at (B) {$B$};
\node[below] at (C) {$C$};
\node[above] at (D) {$D$};

\end{tikzpicture}
\end{center}

\vspace{0.3cm}

\begin{tcolorbox}[
  colback=white,
  colframe=NeutreC!25,
  boxrule=0.4pt,
  height=5.2cm
]
\textbf{Recherche}
\end{tcolorbox}

\end{minipage}

% =========================================================
% PROF -- RENDEZ-VOUS 7
%
% AD = DB = AB
% -> ADB est équilatéral
% -> angle ADB = 60°
%
% A, B, C alignés
% -> angle DBC = 120°
%
% DB = BC
% -> DBC est isocèle en B
%
% Ses deux autres angles valent :
% (180 - 120) / 2 = 30°
%
% ADC = ADB + BDC
%     = 60 + 30
%     = 90°
%
% Donc ACD est rectangle en D.
%
%
% Relances graduées :
%
% 1. "Quelles petites figures reconnais-tu ?"
%
% 2. "Regarde ADB."
%
% 3. "Que sais-tu de ses trois côtés ?"
%
% 4. "Que change l'alignement de A, B et C ?"
%
% 5. "Regarde maintenant DBC."
%
% Pas de démonstration rédigée imposée.
% =========================================================
}

\newpage

% =========================================================
% RENDEZ-VOUS 8
% CONTRÔLER
% =========================================================

\titreRDV
{Rendez-vous 8 -- Où le raisonnement casse-t-il ?}
{Retour sur la situation précédente}

\flashcalcul{
\calcquatre
{\(31{,}2\times2{,}9\approx ?\)}
{\(602\div9\approx ?\)}
{\(25\%\text{ de }96=48\ ?\)}
{\(4{,}08\times24{,}9\approx ?\)}
}

\vspace{0.15cm}

\grandeactivite{A. Trouver la première erreur}{

Voici trois productions concernant la figure du rendez-vous précédent.

\medskip

\textbf{Pour chaque production :}
entoure la première phrase qui n'est pas justifiée.

Si tout est correct, écris \textbf{correct}.

\vspace{0.25cm}

\begin{minipage}[t]{0.315\linewidth}

\begin{tcolorbox}[
  colback=white,
  colframe=NeutreC!45,
  boxrule=0.5pt,
  height=7.2cm
]
\textbf{Production 1}

\medskip

1. \(AD=DB=AB\).

2. Donc \(ADB\) est équilatéral.

3. Donc \(\widehat{ADB}=60^\circ\).

4. \(A,B,C\) sont alignés.

5. Donc \(\widehat{DBC}=120^\circ\).

6. \(DB=BC\), donc \(DBC\) est isocèle en \(B\).

7. Donc \(\widehat{BDC}=30^\circ\).

8. Donc \(\widehat{ADC}=90^\circ\).
\end{tcolorbox}

\end{minipage}
\hfill
\begin{minipage}[t]{0.315\linewidth}

\begin{tcolorbox}[
  colback=white,
  colframe=NeutreC!45,
  boxrule=0.5pt,
  height=7.2cm
]
\textbf{Production 2}

\medskip

1. \(AD=DB=AB\).

2. Donc \(ADB\) est équilatéral.

3. Donc \(\widehat{ADB}=60^\circ\).

4. \(DB=BC\).

5. Donc \(DBC\) est équilatéral.

6. Donc \(\widehat{BDC}=60^\circ\).

7. Donc \(\widehat{ADC}=120^\circ\).
\end{tcolorbox}

\end{minipage}
\hfill
\begin{minipage}[t]{0.315\linewidth}

\begin{tcolorbox}[
  colback=white,
  colframe=NeutreC!45,
  boxrule=0.5pt,
  height=7.2cm
]
\textbf{Production 3}

\medskip

1. \(AD=DB=AB\).

2. Donc \(ADB\) est équilatéral.

3. Donc \(\widehat{ADB}=60^\circ\).

4. \(A,B,C\) sont alignés.

5. Donc \(\widehat{DBC}=120^\circ\).

6. \(DB=BC\), donc \(DBC\) est isocèle en \(B\).

7. Donc \(\widehat{BDC}=60^\circ\).
\end{tcolorbox}

\end{minipage}

% PROF :
% P1 correcte.
% P2 : première erreur = ligne 5.
% P3 : première erreur = ligne 7.
}

\vspace{0.15cm}

\grandeactivite{B. Réparer}{

Choisis une des deux productions incorrectes.

\medskip

\textbf{Change le moins possible pour la rendre correcte.}

\vspace{0.3cm}

\begin{tcolorbox}[
  colback=white,
  colframe=NeutreC!25,
  boxrule=0.4pt,
  height=3.2cm
]
\end{tcolorbox}

% =========================================================
% PROF -- RENDEZ-VOUS 8
%
% Calcul :
%
% 31,2 x 2,9
% proche de 30 x 3 = 90
%
% 602 ÷ 9
% proche de 600 ÷ 9
% soit environ 67
%
% 25 % de 96 = 48 ?
% Faux.
% 25 % = 1/4
% donc 24.
%
% 4,08 x 24,9
% proche de 4 x 25 = 100
%
%
% Nœud :
% ne pas seulement repérer une phrase fausse,
% mais le PREMIER passage qui cesse d'être garanti.
% =========================================================
}

\newpage

% =========================================================
% RENDEZ-VOUS 9
% TRANSFERT / COMMUNICATION
% =========================================================

\titreRDV
{Rendez-vous 9 -- Chercher, organiser, communiquer}
{Situation de transfert}

\flashcalcul{

\begin{tabularx}{\linewidth}{
>{\centering\arraybackslash}X
>{\centering\arraybackslash}X
}
\textbf{Calcul 1} &
\textbf{Calcul 2} \\[0.1cm]

\(28\times25\) &
\(75\%\text{ de }64\) \\[0.2cm]

\multicolumn{1}{c}{
Trouve \textbf{deux façons différentes}.
}
&
\multicolumn{1}{c}{
Trouve \textbf{deux façons différentes}.
}
\\[0.3cm]

\rule{7cm}{0.4pt}
&
\rule{7cm}{0.4pt}
\end{tabularx}
}

\vspace{0.15cm}

\grandeactivite{Une nouvelle figure}{

\begin{minipage}[c]{0.37\linewidth}

\textbf{Le dessin est un croquis.}

\medskip

Les points \(A\), \(B\) et \(C\) sont alignés.

\[
AB=AD
\]

\[
DB=BC
\]

\[
\widehat{BAD}=40^\circ
\]

\vspace{0.45cm}

{\large
\textbf{Détermine la mesure de l'angle \(\widehat{ADC}\).}
}

\vspace{0.6cm}

\begin{tcolorbox}[
  colback=white,
  colframe=NeutreC!25,
  boxrule=0.4pt,
  height=4.8cm
]
\textbf{Recherche}
\end{tcolorbox}

\end{minipage}
\hfill
\begin{minipage}[c]{0.59\linewidth}

\begin{center}
\begin{tikzpicture}[scale=1.4]

% Croquis volontairement non conforme.

\coordinate (A) at (0,0);
\coordinate (B) at (3,0);
\coordinate (C) at (6,0);
\coordinate (D) at (1.35,2.55);

\draw[trait,codeun] (A)--(D);
\draw[trait,codeun] (A)--(B);

\draw[trait,codedeux] (D)--(B);
\draw[trait,codedeux] (B)--(C);

\draw[trait] (D)--(C);

\node[below] at (A) {$A$};
\node[below] at (B) {$B$};
\node[below] at (C) {$C$};
\node[above] at (D) {$D$};

\node at (0.75,0.35) {$40^\circ$};

\end{tikzpicture}
\end{center}

\vspace{0.35cm}

\textbf{Quand ta recherche est terminée :}

\medskip

Explique ton raisonnement en quelques lignes.

\vspace{0.3cm}

\rule{\linewidth}{0.4pt}

\vspace{0.4cm}

\rule{\linewidth}{0.4pt}

\vspace{0.4cm}

\rule{\linewidth}{0.4pt}

\vspace{0.4cm}

\rule{\linewidth}{0.4pt}

\end{minipage}

% =========================================================
% PROF -- RENDEZ-VOUS 9
%
% Calcul 1 :
%
% 28 x 25 = 700
%
% Par exemple :
% 28 x 100 ÷ 4
%
% ou
%
% 28 = 7 x 4
% donc 7 x (4 x 25)
% = 7 x 100
%
%
% Calcul 2 :
%
% 75 % de 64 = 48
%
% par 3/4
%
% ou
%
% 50 % + 25 %
% = 32 + 16
%
%
% GÉOMÉTRIE
%
% AB = AD
% -> ABD isocèle en A.
%
% BAD = 40°
% -> ABD = BDA
% -> chacun vaut 70°.
%
% A, B, C alignés
% -> DBC = 180 - 70 = 110°.
%
% DB = BC
% -> DBC isocèle en B.
%
% BDC = DCB
% -> chacun vaut (180 - 110)/2 = 35°.
%
% ADC = ADB + BDC
%     = 70 + 35
%     = 105°.
%
% Réponse : 105°.
%
%
% Ici seulement, la rédaction devient explicitement
% un objet du travail.
%
% Relances graduées :
%
% 1. "Quelles sous-figures peux-tu étudier ?"
%
% 2. "Commence par ABD."
%
% 3. "Que te donne ensuite l'alignement ?"
%
% 4. "Regarde DBC."
%
% Ne pas transformer ces relances en questions
% imprimées sur la fiche.
% =========================================================
}

\end{document}